\chapter[Sermon 101. Punishment Is Decreed to the Corruptors of This Book. The Lord Says, That He Will Certainly Come to Judgment. The Church Wishes for His Coming.]{Punishment Is Decreed to the Corruptors of This Book. The Lord Says, That He Will Certainly Come to Judgment. The Church Wishes for His Coming.}
\label{sermon:101}
\markright{Sermon 101. Punishment Is Decreed to the Corruptors of This Book ...}

I testify to every person who hears the words of the prophecy of this book: if anyone adds to these things, God will add to him the plagues that are written in this book. And if anyone diminishes from the words of the book of this prophecy, God will take away his part from the book of life, and from the holy city, and from the things which are written in this book. He who testifies these things says, “Surely, it is.” I am coming quickly. Amen. Even so, come, Lord Jesus. The grace of our Lord Jesus Christ be with you all. Amen.

\index{Antichrist}\index{Erasmus of Rotterdam}\index{Jerome, Saint}\index{Tertullian}13. In the concluding section of this decree, a penalty is established for those who despise this book, but especially for counterfeiters or forgers, who (as Dr. Bibliander has said aptly and piously) dare attempt to corrupt or falsify this godly instrument and holy charter of Christ’s empire and bishopric, by adding anything or taking away, or altering the true meaning and sense thereof. This passage is taken outside of the common practice of men. For princes are accustomed at the end of their writings to secure them against detractors by means of threats and warnings. Antichrist, the mimic of our Lord Christ, at the end of his bull adds: “If anyone rashly presumes to go against this our commandment, or presumptuously to infringe the same, let him know that he shall incur the indignation of Almighty God and the blessed Apostles Peter and Paul, and our high displeasure.” And likewise in safeguarding treasures and public things, where danger is feared, they affix writings and seals with wax. For this reason, indeed, knowing that there would be some who would seek to oppress and abolish this book, the Lord sends it well armed to all future generations. We read in ancient authors that certain heretics in the early church greatly asserted themselves in corrupting the scriptures; and that some of them rejected entire books of the holy scripture. And Tertullian attributes the same to Marcion, who also corrupted holy books. Nevertheless, through God’s goodness, it came to pass that we have received the holy books whole and uncorrupted. This is plainly shown by St. Jerome in his commentaries on Isaiah, book 3. And Erasmus of Rotterdam in the Apology of the New Testament, and also in his Apology against James Latomus, and so forth.

However, the Lord at this present time does no new thing, while he commands that nothing should be added or taken away. For once or twice he commanded by Moses: “You shall add nothing to my word, nor take anything from it.” And Solomon in the thirtieth chapter of the Proverbs commands the same. But many marvel and find fault that he has threatened so many plagues to the corrupters. Why then do they not blame and reprove Paul as well, who in one word has encompassed as many plagues and displeasures as John has recited here, where he said to the Galatians: “Although I, or an angel from heaven, should preach to you a gospel other than the one we have preached to you, let him be an outcast or accursed.” And he repeats these same words. Wherefore, if they grant that Paul has offended so little in this regard, that he deserves praise also, let them cease blaming those things which are here most purposefully placed by our Lord Jesus Christ himself through John in their place and time due. Anathema (a word Paul used) is he who is cut off from the fellowship of all good men, devoted to extreme punishment, and even subject to all the evils of both this present life and the life to come. This may be gathered from Deuteronomy and other holy books.

\index{Apocalypse (Book of)}\index{Aquinas, Thomas}And he adds and takes away, not one who uses other diverse and plain words in explaining a sentence of the Apocalypse; but he who puts in anything contrary to the true sense, or varies from those things which are here expressed by the Lord: or he who takes away anything, obscures or perverts what the Lord himself has expressly signified. Thomas of Aquin says he adds, which is to put in a lie; he diminishes, which takes away anything from that which is written therein, or also denies gain, and says the same. Thus much he says. Therefore, this addition and subtraction consists not in words only, but rather in sense. For neither are the prophets, in explaining the law at large, thought to have added anything to God’s word; neither are the Apostles, preaching the liberty of the gospel, said to have taken anything away from the law.

And to testify—that is, to affirm something definite by a declaration or to urge something earnestly, as if to bind the listener—so that he should certainly know that these things spoken have authority over him, and that God will punish unless he obeys.

Concerning the pains or plagues which he threatens at this time, this is spoken of in chapters 15, 16, 17, and 18, and so forth. Likewise, it is declared beforehand that what might be spoken here concerning the book of life and the holy city. Moreover, he comprehends here also all good things in like manner, which are promised in this book to the godly and obedient servants of God, of all which things the contemner, falsifier, and corrupter of this book shall be deprived. With how great evils and dangers than do they entangle themselves, who would have this book utterly suppressed, and prevented from being expounded openly, and coming into the hands of all men? Again, it is most certain that they shall obtain of God all manner of blessing, so many as have a good opinion, and think devoutly of this book, and will set forth and commend unto all men the things that are written in the same to the glory of God, and salvation of the faithful.

\index{John, Saint}14. In the fourteenth place, the authority of this book is sealed, and as it were signed with a subscription. For it follows, he says that bears witness of these things, or he that testifies these things. In general, all the commentators suppose these to be Christ’s words, as though he himself, for a confirmation, had added to it and said: “I Jesus have revealed all these things, and especially those concerning the threatenings against the corrupter, as a true witness, and the same to be undoubted.” For Aquinas says: “Here Christ is brought in, affirming the aforementioned threatening, and approving all things that are written in this book.” But I, for my part, reserving the judgments of others safe, suppose this to be the subscription of Saint John, the writer of this book. For notaries, secretaries, or chancellors of princes are wont at the end of the kings or emperors’ letters or writings to subscribe their name. And indeed Saint John in the History of the Gospel has observed the same manner. For in chapter 19, he says: “And he that saw bore witness, and his testimony is true.” And at the end of the History he subscribes and signs underneath with these words: “This is that disciple which bears witness of these things, and wrote them, and we know that his testimony is true.” For all the church knew and confessed this. In like manner, he seems at this present to have subscribed these things also in his own name and to have said, “He that testifies these things, says.”

15. Next, he brings in again the Lord Jesus himself speaking and promising that he will certainly come for judgment, truly to redeem and glorify the righteous, and to punish the wicked. Therefore, with a strong affirmation, he says, “Even so, surely, and without doubt I am coming.” Although I seem to delay, and to some I do not appear to be coming at all. Nevertheless, I will certainly come in due time, as has been said and declared before. And this same thing is repeated, in similar manner with the same words, more often as a matter most worthy to be noted and known.

He immediately adds the faith and wish, and great desire of Saint John, and of the faithful church, or of any godly person submitting himself to the promise, and saying, “Amen, even so.” That is to say, we acknowledge it to be most certain and undoubted that you promise that you will come. Therefore, we look for you, the Judge of the living and the dead; yea, and we pray with our innermost being, “Come, Lord Jesus.” For elsewhere also we pray daily, “Thy kingdom come.” And all the godly, with sighs unspeakable, wish for the coming of the Judge, for glory, whereof the Apostle treats in the eighth chapter to the Romans, and we have touched the same matter before.

Finally, in his conclusion, he wishes, in the apostolic manner, the grace of our Lord Jesus Christ to all who hear and read this book. Paul writes in his second letter to the Thessalonians, chapter three: “Thus I write in every letter: May the grace of our Lord Jesus Christ be with you all.” Amen. Thus he aligns himself with the apostolic spirit, which is present everywhere. Grace encompasses the entire matter of redemption and the gifts of Christ. Therefore, he wishes to us all the blessing that we have in Christ Jesus our Lord. Of this, the vessel of election has abundantly been endowed. Paul, in chapter one to the Ephesians, prays that the Lord Jesus, who has revealed these holy mysteries to us, may write them in our minds, and deliver us from Antichrist and from all evils; and keep us in the true faith and in his grace. To him be honor and glory, praise and thanks giving, together with the Father and the Holy Spirit, forever: Amen. Come, Lord Jesus, our redeemer and only savior, and glorify those who eagerly await your coming, that we may glorify you forever. Amen.
